\section{Le Méta-Modèle Entity}\label{sub:ent}
Les Entity permettent de simplifier la gestion des données au niveau d'une application, mais aussi de faciliter la sauvegarde en base de données. Plus concrètement, ces Entity nous permettent de prendre en charge la persistance des données de notre application dans une ou plusieurs sources de données, tout en gardant les relations entre celles-ci. Ces composants établissent donc la relation entre notre application et notre bases de données.
Un méta-modèle Entity a été établie par Obeo afin de représenter la structure des Entity qui vont définir la couche métier de notre application. La figure \ref{fig:ent}  montre le méta-modèle Entity simplifié.

\begin{figure}[htb]
  \centering
  \includegraphics[scale=.4]{img/Entity.eps}
  \caption{Méta-modèle Entity}
  \label{fig:ent}
\end{figure}


\subsection{Gestion des entités avec \kwplay{}}
Le framework \kwplay{} utilise la notion de \verb+Model+ qui s'accorde parfaitement avec le concept d'\kwentity{}. 

\subsection{Conception du modèle}
Nous avons employé la simple démarche suivante : pour chaque modèle on associe une entité. À ce stade, on en déduit facilement les attributs de chaque instance d'\verb+Entity+ et les instances \verb+Reference+ correspondant aux associations avec d'autres entités. \kwplay{} utilise des mécanismes d'annotation au sein de ses modèles. Ces annotations permettent notamment de donner des contraintes à des attributs. Nous avons facilement pris en compte ce concept en utilisant la classe \verb+Annotation+ mise à disposition dans les méta datas d'un \verb+Attribut+.
\clearpage


% LocalWords:  Entity framework Model
